
\chapter{振动分析的近似方法}

\intro{精确解析解仅存在于最简单的连续系统中。对于复杂几何、变截面、非均匀材料的结构，必须借助近似方法。本章介绍 Rayleigh 法（能量法）、Rayleigh-Ritz 法、Dunkerley 公式、传递矩阵法和有限元法入门，它们在工程振动分析中广泛应用。}

%%%%%%%%%%%%%%%%%%%%%%%%%%%%%%%%%%%%%%%%%%%%%%%%%%%%%%%%
\section{Rayleigh 法（能量法）}

\subsection{基本原理}

Rayleigh 法（Rayleigh's method）基于机械能守恒原理。对于无阻尼自由振动系统，总能量保持恒定：最大动能等于最大势能。

设梁的振动位移为 $y(x,t) = Y(x)\sin(\omega t + \phi)$，其中 $Y(x)$ 为假设的振型函数。

系统在任意时刻的动能和势能分别为：
\[
T(t) = \frac{1}{2}\int_0^L \rho A(x)\left(\frac{\partial y}{\partial t}\right)^2 dx, \quad
U(t) = \frac{1}{2}\int_0^L EI(x)\left(\frac{\partial^2 y}{\partial x^2}\right)^2 dx
\]

最大动能（速度最大时，$\cos(\omega t+\phi)=1$）：
\[
T_{\max} = \frac{\omega^2}{2}\int_0^L \rho A(x)Y^2(x)\,dx
\]

最大势能（位移最大时，$\sin(\omega t+\phi)=1$）：
\[
U_{\max} = \frac{1}{2}\int_0^L EI(x)[Y''(x)]^2\,dx
\]

由 $T_{\max}=U_{\max}$，得到 \textbf{Rayleigh 商}（Rayleigh quotient）：

\begin{myformula}
\[
\omega^2 = \frac{U_{\max}}{T_{\max}/\omega^2} = \frac{\int_0^L EI(x)[Y''(x)]^2\,dx}{\int_0^L \rho A(x)Y^2(x)\,dx}
\]
\end{myformula}

\begin{mydefinition}{Rayleigh 商的性质}
Rayleigh 商给出的 $\omega^2$ 估计值总是\textbf{大于或等于}真实的基频平方 $\omega_1^2$（上界估计）。当且仅当假设的振型 $Y(x)$ 恰好等于真实的基频振型时，等号成立。
\end{mydefinition}

物理意义：任何不符合真实振型的假设都会引入额外的约束，使系统变``刚''（刚度偏大），从而高估固有频率。因此 Rayleigh 法给出的基频是\textbf{上界}。

\subsection{假设振型的选取原则}

为了获得准确的基频估计，假设的振型函数 $Y(x)$ 应满足：

\begin{enumerate}
    \item \textbf{满足几何边界条件}：挠度 $Y$ 和转角 $Y'$ 的约束必须满足;
    \item \textbf{尽可能满足自然边界条件}：弯矩 $Y''$ 和剪力 $Y'''$ 的条件虽不强制，但满足可提高精度;
    \item \textbf{具有正确的形状}：至少应具有与基频振型相同的凹凸性和节点位置。
\end{enumerate}

一种常用的技巧是\textbf{以静变形曲线作为假设振型}——将结构自身重力或集中力作用于某点产生的静挠度曲线作为 $Y(x)$。这种方法通常给出非常好的基频估计。

\begin{example}
简支等截面梁，弯曲刚度 $EI$，线密度 $\rho A$，长 $L$。

取静变形曲线（均布载荷 $q$ 下的挠度）：
\[
Y(x) = \frac{q}{24EI}(x^4 - 2Lx^3 + L^3 x)
\]
计算 Rayleigh 商：
\[
\omega^2 = \frac{\int_0^L EI[Y''(x)]^2\,dx}{\int_0^L \rho A Y^2(x)\,dx}
    = \frac{EI}{\rho A}\cdot\frac{12\times 31}{L^4\times 1} \cdot ??? 
\]
经过计算可得 $\omega \approx 9.8767\sqrt{EI/(\rho A L^4)}$，而精确解为 $\pi^2\sqrt{EI/(\rho A L^4)} \approx 9.8696\sqrt{EI/(\rho A L^4)}$，误差仅为 $0.07\%$。
\end{example}

另一常见选择：取 $\sin(\pi x/L)$，计算得 $\omega^2 = \pi^4 EI/(\rho A L^4)$，恰好等于精确解——这是因为 $\sin(\pi x/L)$ 本身就是简支梁的精确基频振型。

\begin{figure}[H]\centering
\begin{tikzpicture}[scale=1.3,vib_style/.style={thick},vib_arrow/.style={-Stealth,thick}]
\draw[vib_style] (0,0) -- (5,0) node[right] {$x$};
\draw[vib_style] (0,0) -- (0,2) node[above] {$y$};
\fill (0,0) circle (0.08); \fill (5,0) circle (0.08);
\node at (0,-0.2) {$0$}; \node at (5,-0.2) {$L$};

\draw[domain=0:5,smooth,variable=\x,blue,ultra thick,dashed] plot ({\x},{1.5*(\x*\x*\x*\x/5^4 - 2*\x*\x*\x/5^3 + \x/5)});
\draw[domain=0:5,smooth,variable=\x,red,ultra thick] plot ({\x},{sin(pi*\x/5 r)});

\node[blue] at (2.5,1.0) {静变形假设};
\node[red] at (2.5,0.6) {精确基频振型};
\node at (2.5,-0.7) {Rayleigh 法：假设振型 vs 精确振型};
\end{tikzpicture}
\caption{Rayleigh 法中假设振型（静变形曲线）与精确基频振型的对比。两者形状相似，因此 Rayleigh 法给出高精度的基频估计}
\label{fig:rayleigh-modes}
\end{figure}

\subsection{集中质量与分布质量的混合系统}

Rayleigh 法特别擅长处理既有分布参数又有集中质量的混合系统。假设振型后：

\begin{myformula}
\[
\omega^2 = \frac{U_{\max}}{T^*} = \frac{\int EI[Y''(x)]^2\,dx + \sum k_i Y^2(x_i)}{\int \rho A Y^2(x)\,dx + \sum m_i Y^2(x_i)}
\]
\end{myformula}

其中 $m_i$ 为集中质量，$k_i$ 为集中弹簧刚度，$T^* = T_{\max}/\omega^2$。

这体现了能量法的灵活性——只需分别计算应变能和动能的最大值，即可得到固有频率的估计。

\subsection{多自由度系统的 Rayleigh 商}

对于 $n$ 自由度系统，设振型向量为 $\{\mathbf{Y}\}$，则：

\[
\omega^2 = \frac{\{\mathbf{Y}\}^{\top}\mathbf{K}\{\mathbf{Y}\}}{\{\mathbf{Y}\}^{\top}\mathbf{M}\{\mathbf{Y}\}}
\]

这称为矩阵形式的 Rayleigh 商。任取一个近似的振型向量代入，即可得到基频的上界估计。

--------------------------------------------------------------------
\section{Rayleigh-Ritz 法}

\subsection{基本思想}

Rayleigh 法仅能得到基频的估计（且只能估计上界）。\textbf{Rayleigh-Ritz 法}将 Rayleigh 法推广：取一组线性无关的基函数，将假设的振型表示为这些基函数的线性组合，通过变分原理将泛函极值问题转化为代数特征值问题。这样可以同时估计多阶固有频率。

\subsection{数学推导}

假设振型可以近似为 $N$ 个基函数 $\psi_i(x)$ 的线性组合：
\[
Y(x) = \sum_{i=1}^{N} a_i \psi_i(x) = \{\bm{\psi}(x)\}^{\top}\{\mathbf{a}\}
\]

其中 $\{\mathbf{a}\} = \{a_1, a_2, \dots, a_N\}^{\top}$ 为待定系数，$\{\bm{\psi}\}$ 为基函数向量。每个基函数 $\psi_i(x)$ 必须满足几何边界条件。

将 $Y(x)$ 代入 Rayleigh 商：
\[
\omega^2 = \frac{\{\mathbf{a}\}^{\top}\mathbf{K}\{\mathbf{a}\}}{\{\mathbf{a}\}^{\top}\mathbf{M}\{\mathbf{a}\}}
\]

其中：
\[
K_{ij} = \int_0^L EI(x)\,\psi_i''(x)\psi_j''(x)\,dx, \quad
M_{ij} = \int_0^L \rho A(x)\,\psi_i(x)\psi_j(x)\,dx
\]

$\mathbf{K}$ 和 $\mathbf{M}$ 为 $N\times N$ 的对称矩阵，分别称为\textbf{广义刚度矩阵}和\textbf{广义质量矩阵}。

Rayleigh 商取驻值的条件为 $\partial(\omega^2)/\partial a_i = 0$（$i=1,\dots,N$），这等价于广义特征值问题：

\begin{myformula}
\[
\mathbf{K}\{\mathbf{a}\} = \omega^2 \mathbf{M}\{\mathbf{a}\}
\]
\end{myformula}

求解此 $N$ 维代数特征值问题，得到 $N$ 个近似的固有频率 $\omega_1 \le \omega_2 \le \dots \le \omega_N$ 和对应的特征向量 $\{\mathbf{a}^{(k)}\}$。将特征向量代入基函数组合即得近似的固有振型。

\begin{theorem}[Rayleigh-Ritz 的收敛性]
Rayleigh-Ritz 法给出的近似固有频率是真实固有频率的上界：
\[
\omega_k^{\text{(Ritz)}} \ge \omega_k^{\text{(true)}}, \quad k=1,2,\dots,N
\]
随着基函数个数 $N$ 的增加，近似值从上方单调收敛到真实值。
\end{theorem}

\subsection{基函数的选取}

常用的基函数类型包括：

\begin{enumerate}
    \item \textbf{幂函数}：$\psi_i(x) = x^{i-1}$（仅当满足几何边界条件时适用，如悬臂梁取 $\psi_i(x)=x^{i+1}$）;
    \item \textbf{三角函数}：$\psi_i(x) = \sin(i\pi x/L)$——自然满足简支边界条件;
    \item \textbf{多项式}：$\psi_i(x) = (x/L)^{i-1}\cdot\text{(几何边界因子)}$——如悬臂梁取 $(x/L)^{i+1}$;
    \item \textbf{梁函数}：以均匀梁的精确振型为基函数——这是最高效的选择。
\end{enumerate}

\begin{remark}
基函数只需满足几何边界条件。满足自然边界条件的基函数可提高收敛速度，但不是必须的。这是 Rayleigh-Ritz 法相对于直接求解 PDE 的一大优势。
\end{remark}

\begin{figure}[H]\centering
\begin{tikzpicture}[scale=1.1,vib_style/.style={thick},vib_arrow/.style={-Stealth,thick}]
\draw[vib_style] (0,0) -- (0,3.5) node[above] {$\omega$};
\draw[vib_style] (0,0) -- (6,0) node[right] {$N$（基函数个数）};
\draw[vib_style,dashed,gray] (0,2.5) -- (6,2.5) node[right] {$\omega_1^{\text{true}}$};
\draw[vib_style,dashed,gray] (0,3.2) -- (6,3.2) node[right] {$\omega_2^{\text{true}}$};

\fill (1,2.7) circle (0.08); \fill (2,2.58) circle (0.08);
\fill (3,2.53) circle (0.08); \fill (4,2.51) circle (0.08);
\fill (5,2.505) circle (0.08);
\draw[blue,ultra thick] (1,2.7) -- (2,2.58) -- (3,2.53) -- (4,2.51) -- (5,2.505);
\node[blue] at (3,2.8) {$\omega_1^{\text{(Ritz)}}$（从上方收敛）};

\fill (1,3.4) circle (0.08); \fill (2,3.32) circle (0.08);
\fill (3,3.26) circle (0.08); \fill (4,3.22) circle (0.08);
\fill (5,3.21) circle (0.08);
\draw[red,ultra thick] (1,3.4) -- (2,3.32) -- (3,3.26) -- (4,3.22) -- (5,3.21);
\node[red] at (3,3.45) {$\omega_2^{\text{(Ritz)}}$（从上方收敛）};
\end{tikzpicture}
\caption{Rayleigh-Ritz 法的收敛性质。近似固有频率从真实值的上方单调收敛}
\label{fig:ritz-convergence}
\end{figure}

--------------------------------------------------------------------
\section{Dunkerley 公式}

\subsection{基本思想}

Rayleigh 法给出基频的上界估计。Dunkerley 公式则从另一个角度出发，给出基频的\textbf{下界估计}。两者配合使用可以将基频的真实值限制在一个区间内。

\subsection{公式推导}

考虑一个多自由度系统，其柔度矩阵为 $\mathbf{A} = \mathbf{K}^{-1}$。Dunkerley 基于如下观察：系统的总柔度可以近似为各元件柔度之和。

对于轴上有多个圆盘的系统（或梁上有多个集中质量），Dunkerley 公式为：

\begin{myformula}
\[
\frac{1}{\omega_1^2} \approx \sum_{i=1}^{n} \frac{1}{\omega_{ii}^2}
\]
\end{myformula}

其中 $\omega_{ii}$ 是系统中仅保留第 $i$ 个质量（其他质量设为零）时的基频。对于梁：
\[
\omega_{ii}^2 = \frac{1}{m_i \alpha_{ii}}, \quad \alpha_{ii} = \text{在位置 } i \text{ 处施加单位力产生的位移}
\]

Dunkerley 公式也可推广到包含分布质量的系统：
\[
\frac{1}{\omega_1^2} \approx \frac{1}{\omega_{d}^2} + \sum_{i}\frac{1}{\omega_{ii}^2}
\]
其中 $\omega_d$ 是仅考虑分布质量时的基频。

\begin{remark}
Dunkerley 公式的关键性质：它给出的 $\omega_1$ 估计值总是\textbf{小于或等于}真实的基频（下界）。这正好与 Rayleigh 法互补：
\[
\omega_1^{\text{(Dunkerley)}} \le \omega_1^{\text{(true)}} \le \omega_1^{\text{(Rayleigh)}}
\]
\end{remark}

\begin{example}
简支梁中点处有集中质量 $M$。梁单位长度质量 $\rho A$，弯曲刚度 $EI$。

仅考虑分布质量：$\omega_d = \pi^2\sqrt{EI/(\rho A L^4)}$。

仅考虑集中质量：$\omega_{11} = \sqrt{48EI/(M L^3)}$。

Dunkerley 下界：
\[
\frac{1}{\omega_1^2} \approx \frac{1}{\omega_d^2} + \frac{1}{\omega_{11}^2}
\]
\end{example}

--------------------------------------------------------------------
\section{传递矩阵法（Transfer Matrix Method）}

\subsection{基本概念}

传递矩阵法（TMM）是分析链式结构（如多跨轴、多盘转子、阶梯梁）振动的有力工具。其核心思想是：将结构分解为一系列单元，每个单元的状态向量（位移、转角、弯矩、剪力）通过传递矩阵从前一单元的右端传递到下一单元的右端。

\begin{mydefinition}{状态向量}
梁单元的\textbf{状态向量}定义为四个力学量组成的列向量：
\[
\{\mathbf{Z}(x)\} = \{y(x),\; \theta(x),\; M(x),\; V(x)\}^{\top}
\]
\end{mydefinition}

状态向量的``正方向''约定：
\[
\theta = \frac{dy}{dx}, \quad M = EI\frac{d^2y}{dx^2}, \quad V = EI\frac{d^3y}{dx^3}
\]

Sign convention: displacement $y$ upwards positive, slope $\theta$ counterclockwise positive.

\subsection{场矩阵——穿过均匀场段}

对于长度为 $\ell$ 的均匀梁段（场单元），$i$ 端（左端）与 $i+1$ 端（右端）的状态向量之间有关系：
\[
\{\mathbf{Z}\}_{i+1} = \mathbf{F}\{\mathbf{Z}\}_i
\]

其中 $\mathbf{F}$ 为\textbf{场矩阵}（field matrix）。对 Euler-Bernoulli 梁的无质量场段：

\begin{myformula}
\[
\mathbf{F} = \begin{bmatrix}
1 & \ell & \frac{\ell^2}{2EI} & \frac{\ell^3}{6EI} \\[4pt]
0 & 1 & \frac{\ell}{EI} & \frac{\ell^2}{2EI} \\[4pt]
0 & 0 & 1 & \ell \\[4pt]
0 & 0 & 0 & 1
\end{bmatrix}
\]
\end{myformula}

如果考虑分布质量，场矩阵将涉及 $\sin$, $\cos$, $\sinh$, $\cosh$ 函数（包含振动频率 $\omega$）。

\subsection{点矩阵——跨越集中元件}

跨越集中质量 $m$（或弹簧 $k$、刚性支承）时，状态向量发生跳跃。以集中质量为例：

\[
\{\mathbf{Z}\}_{\text{right}} = \mathbf{P}\{\mathbf{Z}\}_{\text{left}}
\]

\begin{myformula}
\[
\mathbf{P}_{\text{mass}} = \begin{bmatrix}
1 & 0 & 0 & 0 \\
0 & 1 & 0 & 0 \\
0 & 0 & 1 & 0 \\
m\omega^2 & 0 & 0 & 1
\end{bmatrix}, \quad
\mathbf{P}_{\text{spring}} = \begin{bmatrix}
1 & 0 & 0 & 0 \\
0 & 1 & 0 & 0 \\
0 & 0 & 1 & 0 \\
-k & 0 & 0 & 1
\end{bmatrix}
\]
\end{myformula}

\subsection{整体传递矩阵}

对于由 $n$ 个单元组成的链式结构，从最左端到最右端的整体传递关系为：
\[
\{\mathbf{Z}\}_n = \mathbf{T}\,\{\mathbf{Z}\}_0
\]
\[
\mathbf{T} = \mathbf{F}_n\mathbf{P}_n\mathbf{F}_{n-1}\cdots\mathbf{P}_2\mathbf{F}_1
\]

其中 $\mathbf{F}_i$ 为第 $i$ 个场单元的场矩阵，$\mathbf{P}_i$ 为第 $i$ 个点单元的矩阵。

引入边界条件后，总可写为 $\mathbf{T}_{BC}(\omega)\{\mathbf{Z}\}_0 = \{\mathbf{0}\}$。非零解的条件是系数行列式为零：

\[
\det\mathbf{T}_{BC}(\omega) = 0
\]

这就是 TMM 的\textbf{频率方程}。通过数值搜索满足此方程的 $\omega$ 值即可获得固有频率。

\begin{remark}
TMM 的优点：(1) 矩阵规模不随单元数增长——始终是 $4\times 4$;(2) 适用于任意复杂的链式结构;(3) 编程实现简单。缺点：对长链结构可能出现数值不稳定（需要改进的 Riccati 传递矩阵法）。
\end{remark}

\begin{figure}[H]\centering
\begin{tikzpicture}[scale=1.0,vib_style/.style={thick},vib_arrow/.style={-Stealth,thick}]
\draw[vib_style,fill=gray!20] (0,0) rectangle (1.5,1.5);
\draw[vib_style,fill=gray!20] (3,0) rectangle (4.5,1.5);
\draw[vib_style,fill=gray!20] (6,0) rectangle (7.5,1.5);
\fill (1.6,0.75) circle (0.25) node {$m_1$};
\fill (4.6,0.75) circle (0.25) node {$m_2$};
\node at (0.75,0.3) {$\ell_1$};
\node at (3.75,0.3) {$\ell_2$};
\node at (6.75,0.3) {$\ell_3$};
\node at (0.75,-0.4) {场段1};
\node at (3.75,-0.4) {场段2};
\node at (6.75,-0.4) {场段3};

\draw[vib_arrow] (1.75,2.2) -- (1.75,1.25) node[midway,right] {$\mathbf{P}_1$};
\draw[vib_arrow] (4.75,2.2) -- (4.75,1.25) node[midway,right] {$\mathbf{P}_2$};
\draw[vib_arrow] (1,2.2) -- (1,1.7) node[midway,left] {$\mathbf{F}_1$};
\draw[vib_arrow] (4,2.2) -- (4,1.7) node[midway,left] {$\mathbf{F}_2$};
\draw[vib_arrow] (7,2.2) -- (7,1.7) node[midway,left] {$\mathbf{F}_3$};

\node at (-0.5,0.75) {$\{\mathbf{Z}\}_0$};
\node at (8.0,0.75) {$\{\mathbf{Z}\}_3$};

\draw[vib_arrow] (-0.2,0.75) -- (0,0.75);
\draw[vib_arrow] (7.5,0.75) -- (7.7,0.75);
\end{tikzpicture}
\caption{传递矩阵法的链式结构模型。场段通过场矩阵 $\mathbf{F}$ 传递，集中元件通过点矩阵 $\mathbf{P}$ 连接}
\label{fig:transfer-matrix-model}
\end{figure}

--------------------------------------------------------------------
\section{有限元法入门}

\subsection{基本思想}

有限元法（FEM）将连续结构离散为有限个\textbf{单元}（elements），单元之间通过\textbf{节点}（nodes）连接。在每个单元内假设位移插值函数（形函数），通过能量原理建立单元方程，再组装成整体方程求解。

FEM 求解振动问题的步骤：
\begin{enumerate}
    \item \textbf{离散化}：将结构划分为有限个单元，定义节点;
    \item \textbf{建立单元方程}：对每个单元计算单元刚度矩阵 $\mathbf{k}^{(e)}$ 和单元质量矩阵 $\mathbf{m}^{(e)}$;
    \item \textbf{组装}：将单元矩阵组装成整体刚度矩阵 $\mathbf{K}$ 和整体质量矩阵 $\mathbf{M}$;
    \item \textbf{施加边界条件}：消去约束自由度;
    \item \textbf{求解特征值问题}：$(\mathbf{K} - \omega^2\mathbf{M})\{\mathbf{X}\} = \{\mathbf{0}\}$。
\end{enumerate}

\subsection{梁单元的刚度矩阵和质量矩阵}

对于平面二节点梁单元，每节点两个自由度（挠度 $v$ 和转角 $\theta$），单元有四个自由度：
\[
\{\mathbf{d}\}^{(e)} = \{v_1,\; \theta_1,\; v_2,\; \theta_2\}^{\top}
\]

位移插值采用三次 Hermite 多项式：
\[
v(\xi) = N_1 v_1 + N_2\theta_1 + N_3 v_2 + N_4\theta_2, \quad \xi = x/\ell \in [0,1]
\]
\[
N_1 = 1-3\xi^2+2\xi^3,\; N_2 = \ell(\xi-2\xi^2+\xi^3),\; N_3 = 3\xi^2-2\xi^3,\; N_4 = \ell(-\xi^2+\xi^3)
\]

\textbf{单元刚度矩阵}（弯曲）：
\begin{myformula}
\[
\mathbf{k}^{(e)} = \frac{EI}{\ell^3}\begin{bmatrix}
12 & 6\ell & -12 & 6\ell \\
6\ell & 4\ell^2 & -6\ell & 2\ell^2 \\
-12 & -6\ell & 12 & -6\ell \\
6\ell & 2\ell^2 & -6\ell & 4\ell^2
\end{bmatrix}
\]
\end{myformula}

\textbf{一致质量矩阵}（consistent mass matrix，由动能积分导出）：
\begin{myformula}
\[
\mathbf{m}^{(e)} = \frac{\rho A \ell}{420}\begin{bmatrix}
156 & 22\ell & 54 & -13\ell \\
22\ell & 4\ell^2 & 13\ell & -3\ell^2 \\
54 & 13\ell & 156 & -22\ell \\
-13\ell & -3\ell^2 & -22\ell & 4\ell^2
\end{bmatrix}
\]
\end{myformula}

工程中也常用\textbf{集中质量矩阵}（lumped mass matrix，将质量集中到节点上）：
\[
\mathbf{m}_{\text{lumped}}^{(e)} = \frac{\rho A\ell}{2}\operatorname{diag}(1, 0, 1, 0)
\]

\begin{remark}
一致质量矩阵给出的固有频率是精确解的上界（与 Rayleigh-Ritz 类似），而集中质量矩阵则通常给出偏低的估计。两种质量矩阵的组合在工程实践中各有适用场景。
\end{remark}

\subsection{组装与全局特征值问题}

将各单元的刚度矩阵和质量矩阵按节点自由度编号``装配''到全局矩阵中。设节点自由度总数为 $N$，则全局特征值问题为：

\[
(\mathbf{K}_{N\times N} - \omega^2\mathbf{M}_{N\times N})\{\mathbf{X}\} = \{\mathbf{0}\}
\]

施加边界条件后（消去约束自由度），实际求解的是缩减后的特征值问题，维度等于系统的有效自由度数。

\subsection{FEM 与其他方法的比较}

\begin{table}[H]\centering
\caption{振动分析方法比较}
\label{tab:method-comparison}
\begin{tabular}{lccc}
\toprule
方法 & 精度 & 适用范围 & 计算量 \\
\midrule
精确解法 & 精确 & 简单均匀结构 & — \\
Rayleigh 法 & 基频上界 & 任意结构 & 极小 \\
Rayleigh-Ritz & 可多阶，均上界 & 任意结构 & 中等 \\
Dunkerley & 基频下界 & 轴系、集中质量 & 极小 \\
传递矩阵法 & 中等 & 链式结构 & 中等 \\
有限元法 & 可任意精度 & 任意结构 & 大 \\
\bottomrule
\end{tabular}
\end{table}

--------------------------------------------------------------------
\section{数值模拟与代码实现}

\begin{mycode}{Rayleigh-Ritz 法求解简支梁的固有频率}
\begin{lstlisting}[language=Python]
import numpy as np
from scipy.linalg import eigh
import matplotlib.pyplot as plt

L = 1.0; EI = 1.0; rhoA = 1.0

def psi(i, x):
    return np.sin((i + 1) * np.pi * x / L)

def psi_dd(i, x):
    n = i + 1
    return -(n * np.pi / L)**2 * np.sin(n * np.pi * x / L)

N = 6  # 基函数个数
n_gauss = 20
x_gauss = np.linspace(0, L, n_gauss)
dx = L / (n_gauss - 1)

K = np.zeros((N, N))
M = np.zeros((N, N))

for i in range(N):
    for j in range(N):
        K[i, j] = np.trapz(EI * psi_dd(i, x_gauss) * psi_dd(j, x_gauss), x_gauss)
        M[i, j] = np.trapz(rhoA * psi(i, x_gauss) * psi(j, x_gauss), x_gauss)

eigenvalues, eigenvectors = eigh(K, M)
omega_ritz = np.sqrt(eigenvalues)
omega_exact = np.array([(n * np.pi / L)**2 * np.sqrt(EI / rhoA) for n in range(1, N+1)])

print("Mode | Ritz-Ritz     | Exact         | Error (%)")
print("-" * 55)
for i in range(N):
    err = abs(omega_ritz[i] - omega_exact[i]) / omega_exact[i] * 100
    print(f"  {i+1}  | {omega_ritz[i]:.6f}    | {omega_exact[i]:.6f}    | {err:.4f}")
\end{lstlisting}
\end{mycode}

\begin{mycode}{传递矩阵法求解多盘转子的临界转速}
\begin{lstlisting}[language=Python]
import numpy as np
from scipy.optimize import brentq

EI = 1e5; L = 1.0; nelem = 5
le = L / nelem
masses = np.array([0.5, 1.0, 1.5, 1.0, 0.5])

def field_matrix(omega, le, EI):
    F = np.eye(4)
    F[0, 1] = le; F[0, 2] = le**2 / (2 * EI); F[0, 3] = le**3 / (6 * EI)
    F[1, 2] = le / EI; F[1, 3] = le**2 / (2 * EI)
    F[2, 3] = le
    return F

def point_matrix(m, omega):
    P = np.eye(4)
    P[3, 0] = m * omega**2
    return P

def overall_matrix(omega):
    T = np.eye(4)
    for i in range(nelem):
        T = field_matrix(omega, le, EI) @ T
        if i < nelem:
            T = point_matrix(masses[i], omega) @ T
    return T

def det_BC(omega):
    T = overall_matrix(omega)
    return T[1, 1] * T[2, 3] - T[1, 3] * T[2, 1]

omega_guesses = np.linspace(100, 5000, 50)
freqs = []
for i in range(len(omega_guesses) - 1):
    a, b = omega_guesses[i], omega_guesses[i+1]
    if det_BC(a) * det_BC(b) < 0:
        root = brentq(det_BC, a, b, xtol=1e-8)
        freqs.append(root)

print("临界转速 (rad/s):", np.round(freqs, 2))
\end{lstlisting}
\end{mycode}

\begin{mycode}{简单有限元法求解梁的固有频率}
\begin{lstlisting}[language=Python]
import numpy as np
from scipy.linalg import eigh

def beam_element_matrices(EI, rhoA, L):
    ke = EI / L**3 * np.array([
        [12, 6*L, -12, 6*L],
        [6*L, 4*L**2, -6*L, 2*L**2],
        [-12, -6*L, 12, -6*L],
        [6*L, 2*L**2, -6*L, 4*L**2]
    ])
    me_lumped = rhoA * L / 2 * np.diag([1, 0, 1, 0])
    me_consistent = rhoA * L / 420 * np.array([
        [156, 22*L, 54, -13*L],
        [22*L, 4*L**2, 13*L, -3*L**2],
        [54, 13*L, 156, -22*L],
        [-13*L, -3*L**2, -22*L, 4*L**2]
    ])
    return ke, me_lumped, me_consistent

EI = 1.0; rhoA = 1.0; L_total = 1.0
n_elements = 10
le = L_total / n_elements
n_dof = 2 * (n_elements + 1)

K_global = np.zeros((n_dof, n_dof))
M_global = np.zeros((n_dof, n_dof))

ke, me_l, me_c = beam_element_matrices(EI, rhoA, le)
me = me_l  # 使用集中质量矩阵

for e in range(n_elements):
    dof = [2*e, 2*e+1, 2*e+2, 2*e+3]
    for i in range(4):
        for j in range(4):
            K_global[dof[i], dof[j]] += ke[i, j]
            M_global[dof[i], dof[j]] += me[i, j]

# 简支边界条件：消去自由度 0（挠度）和 2*n_elements（挠度）
fixed_dofs = [0, 2 * n_elements]
free_dofs = [i for i in range(n_dof) if i not in fixed_dofs]

K_red = K_global[np.ix_(free_dofs, free_dofs)]
M_red = M_global[np.ix_(free_dofs, free_dofs)]

eigvals, eigvecs = eigh(K_red, M_red)
omega_fem = np.sqrt(eigvals[:5])
omega_exact = np.array([(n * np.pi / L_total)**2 * np.sqrt(EI / rhoA) for n in range(1, 6)])

print("Mode | FEM (Hz)       | Exact (Hz)     | Error (%)")
print("-" * 55)
for i in range(5):
    f_fem = omega_fem[i] / (2 * np.pi)
    f_exact = omega_exact[i] / (2 * np.pi)
    err = abs(f_fem - f_exact) / f_exact * 100
    print(f"  {i+1}  | {f_fem:.6f}      | {f_exact:.6f}      | {err:.4f}")
\end{lstlisting}
\end{mycode}

--------------------------------------------------------------------
\section{小结}

本章介绍的五种近似方法各有特色和适用范围：

\begin{enumerate}
    \item \textbf{Rayleigh 法}：基于能量守恒，给出基频的\textbf{上界}估计。操作简单，适合快速估算;
    \item \textbf{Rayleigh-Ritz 法}：将 Rayleigh 法推广到多自由度，可同时估计多阶频率（均为上界），随着基函数增多而收敛;
    \item \textbf{Dunkerley 公式}：给出基频的\textbf{下界}估计，与 Rayleigh 配合可确定频率范围;
    \item \textbf{传递矩阵法}：适合链式结构，矩阵规模恒定（$4\times4$），频率方程为一维数值搜索;
    \item \textbf{有限元法}：最通用，可处理任意复杂结构，精度随网格细化提高，是现代工程振动分析的标准工具。
\end{enumerate}

这些方法的共同思路是：用有限的计算资源获得满足工程精度要求的振动特性估计，体现了分析力学与数值计算的完美结合。

